%% LyX 2.4.4 created this file.  For more info, see https://www.lyx.org/.
%% Do not edit unless you really know what you are doing.
\documentclass[english]{article}
\usepackage[T1]{fontenc}
\usepackage[latin9]{inputenc}
\usepackage{enumitem}
\usepackage{amsmath}
\usepackage{amsthm}
\usepackage{amssymb}
\usepackage{geometry}
\geometry{verbose,tmargin=1in,bmargin=1in,lmargin=1in,rmargin=1in}
\PassOptionsToPackage{normalem}{ulem}
\usepackage{ulem}

\makeatletter
%%%%%%%%%%%%%%%%%%%%%%%%%%%%%% Textclass specific LaTeX commands.
\numberwithin{equation}{section}
\numberwithin{figure}{section}
\newlength{\lyxlabelwidth}      % auxiliary length 
\theoremstyle{plain}
\newtheorem*{prop*}{\protect\propositionname}
\theoremstyle{definition}
\newtheorem*{example*}{\protect\examplename}
\theoremstyle{plain}
\newtheorem*{thm*}{\protect\theoremname}
\theoremstyle{definition}
\newtheorem*{defn*}{\protect\definitionname}

%%%%%%%%%%%%%%%%%%%%%%%%%%%%%% User specified LaTeX commands.
\usepackage{fancyhdr}
\usepackage{lastpage}
\pagestyle{fancy}
\usepackage{enumitem}
\usepackage{ifthen}
\everymath=\expandafter{\the\everymath\displaystyle}
%\DeclareMathSizes{10}{10}{10}{10}
\usepackage{multicol}

\usepackage{tikz}
\usepackage{tikz-3dplot}
\usetikzlibrary{calc,arrows}

\usetikzlibrary{patterns}
\usetikzlibrary{plotmarks}
\usepackage{pgfplots}
\pgfplotsset{compat=newest}

\usepackage{calc}
\usepackage[nomessages]{fp}

\usepackage{datetime}
% Added by lyx2lyx
\setlength{\parskip}{\smallskipamount}
\setlength{\parindent}{0pt}

\makeatother

\usepackage{babel}
\providecommand{\definitionname}{Definition}
\providecommand{\examplename}{Example}
\providecommand{\propositionname}{Proposition}
\providecommand{\theoremname}{Theorem}

\begin{document}
\renewcommand{\author}{Dr.\,James Ashe}

\newcommand{\classNum}{336}

\newcommand{\classSec}{A}

\newcommand{\examType}{Lecture Notes 3.4}

\renewcommand{\date}{\formatdate{18}{10}{2013}}

%%First page's header%%%%%%%%%%%%%%%%%%%%%%
\fancypagestyle{firststyle} %%header settings for first pagr
{
	\fancyhead[L]{MTH \classNum{}(\classSec{}) \author{}\\
	\textbf{\examType{}} ~\date{}}
	\fancyhead[C]{}
	\fancyhead[R]{\textbf{}}
	\setlength{\headsep}{14pt}
}
%%%%%%%%%%%%%%%%%%%%%%%%%%%%%%%%%%%%%%%%%%

%%Next page's header%%%%%%%%%%%%%%%%%%%%%%%
\setlist{nolistsep}
\lhead{\classNum{}(\classSec{})} 
\chead{\textbf{}} 
\cfoot{\thepage{}~of~\pageref{LastPage}} 
\setlength{\headsep}{5pt} 
%%%%%%%%%%%%%%%%%%%%%%%%%%%%%%%%%%%%%%%%%%%

%%Various settings; see the preamble for other settings%%%%%
%\setenumerate[0]{leftmargin=12pt,itemindent=0pt,labelwidth=0pt}
\setlist{topsep=.5pt}
\renewcommand{\labelenumi}{\textbf{\arabic{enumi}.}}
\renewcommand{\labelenumii}{\textbf{\alph{enumii}.}}
\thispagestyle{firststyle}
%%%%%%%%%%%%%%%%%%%%%%%%%%%%%%%%%%%%%%%%%%

\newcommand{\xmax}{0}
\newcommand{\xmin}{-\xmax}
\newcommand{\xstep}{0}
\newcommand{\ymax}{0}
\newcommand{\ymin}{-\ymax}
\newcommand{\ystep}{0} 
\newcommand{\dspace}{\vspace{1cm}}
\newcommand{\xa}{0}
\newcommand{\xb}{0}
\newcommand{\ya}{1}
\newcommand{\yb}{1}

\subsubsection*{3.3 Review: Linear Independence and Basis}
\begin{prop*}
\textbf{\textup{(3.4) }}\textup{\emph{Let $A\in\mathcal{M}_{nn}$.
Then $A$ is nonsingular if and only if its column vectors form a
basis for $\mathbb{R}^{n}$.}}
\end{prop*}
\begin{example*}
Show that $\mathbf{v_{1}}=(1,0,3)$, $\mathbf{v}_{2}=(1,1,2)$, $\mathbf{v}_{3}=(1,1,1)$,
and $\mathbf{b}=(3,0,1)$. Show that $\{\mathbf{v_{1}},\mathbf{v_{2}},\mathbf{v}_{3}\}$
is a basis for $\mathbb{R}^{3}$.
\[
A=\begin{bmatrix}\begin{array}{rrr}
1 & 1 & 1\\
0 & 1 & 1\\
3 & 2 & 1
\end{array}\end{bmatrix}\rightsquigarrow\begin{bmatrix}\begin{array}{rrr}
1 & 0 & 0\\
0 & 1 & 0\\
0 & 0 & 1
\end{array}\end{bmatrix}
\]
So it is a basis by the proposition. 
\end{example*}
\begin{itemize}
\item $A\mathbf{x}=\mathbf{0}$ has only the trivial solution$\Rightarrow$the
vectors are linearly indepedent.
\item $A\mathbf{x}=\mathbf{b}$ has a unique solution$\Rightarrow$any vector
in $\mathbb{R}^{3}$ is a linear combination of the vectors . 
\item Together this means the vectors form a basis of $\mathbb{R}^{3}$.
\end{itemize}
To write $\mathbf{b}$ as a linear combination of $\{\mathbf{v_{1}},\mathbf{v_{2}},\mathbf{v}_{3}\}$,
reduce
\[
\begin{bmatrix}\begin{array}{rrr|}
1 & 1 & 1\\
0 & 1 & 1\\
3 & 2 & 1
\end{array}\begin{array}{r}
3\\
0\\
1
\end{array}\end{bmatrix}\rightsquigarrow\begin{bmatrix}\begin{array}{rrr|}
1 & 0 & 0\\
0 & 1 & 0\\
0 & 0 & 1
\end{array}\begin{array}{r}
3\\
-8\\
8
\end{array}\end{bmatrix}\mbox{.}
\]
Which gives the \emph{coordinates }of $\mathbf{b}$ with respect to
$\{\mathbf{v_{1}},\mathbf{v_{2}},\mathbf{v}_{3}\}$,
\[
3\begin{bmatrix}\begin{array}{r}
1\\
0\\
3
\end{array}\end{bmatrix}-8\begin{bmatrix}\begin{array}{r}
1\\
1\\
2
\end{array}\end{bmatrix}+8\begin{bmatrix}\begin{array}{r}
1\\
1\\
1
\end{array}\end{bmatrix}=\begin{bmatrix}\begin{array}{r}
3\\
0\\
1
\end{array}\end{bmatrix}
\]


\subsubsection*{3.4 Dimensions and Its Consequences. }
\begin{itemize}
\item There are many bases for a given subspace $V\subset\mathbb{R}^{n}$,
but they ALL have the same number of elements.
\end{itemize}
\begin{prop*}
Let $V\subset\mathbb{R}^{n}$ be a subspace, let $\{\mathbf{v}_{1},\dots,\mathbf{v}_{k}\}$
be a basis for $V$, and let $\{\mathbf{w}_{1},\dots,\mathbf{w}_{\ell}\}\in V$.
If $\ell>k$ then $\{\mathbf{w}_{1},\dots,\mathbf{w}_{\ell}\}$ must
be \uline{linearly dependent}.
\end{prop*}
\begin{example*}
Let $V=\{\mathbf{v}_{1},\dots,\mathbf{v}_{n+1}\}$ for $v_{i}\in\mathbb{R}^{n}$. 
\end{example*}
\begin{enumerate}
\item Is $V$ a basis for $\mathbb{R}^{n}$? 
\item Is $V$ linearly independent or dependent?
\end{enumerate}
\begin{thm*}
Let $V\subset\mathbb{R}^{n}$ be a subspace, and let $\{\mathbf{v}_{1},\dots,\mathbf{v}_{k}\}$
and\textup{ $\{\mathbf{w}_{1},\dots,\mathbf{w}_{\ell}\}$ }\textup{\emph{be
two bases for $V$, then $k=\ell$.}}
\end{thm*}
\begin{itemize}
\item Recall that the standard basis for $\mathbb{R}^{n}$ has $n$ vectors
so any other basis of $\mathbb{R}^{n}$ will also have $n$ vectors.
\end{itemize}
\begin{defn*}
The \emph{dimension }is a subspace $V\subset\mathbb{R}^{n}$ is the
number of vectors in any basis for $V$, denoted $\dim V$.
\end{defn*}
\begin{itemize}
\item By convecntion $\dim\{\mathbf{0}\}=0$.
\item $\dim\mathbb{R}^{n}=n$ ($\mathbb{R}^{n}$ is a subspace of itself).
\end{itemize}
\begin{example*}
(HW 1.b) Find a basis for each of the given subspaces and determine
its dimension.
\end{example*}
\begin{enumerate}
\item $V=\left\{ \mathbf{x}\in\mathbb{R}^{n}\,:\,x_{1}+x_{2}+x_{3}+x_{4}=0,\,x_{2}+x_{4}=0\right\} \subset\mathbb{R}^{4}$.
\[
\begin{bmatrix}\begin{array}{rrrr|r}
1 & 1 & 1 & 1 & 0\\
0 & 1 & 0 & 1 & 0
\end{array}\end{bmatrix}\rightsquigarrow\begin{bmatrix}\begin{array}{rrrr|r}
1 & 0 & 1 & 0 & 0\\
0 & 1 & 0 & 1 & 0
\end{array}\end{bmatrix}\Rightarrow x_{1}=-x_{3}\mbox{ and }x_{2}=-x_{4}
\]
$\Rightarrow$the general solution is $\mathbf{x}=x_{3}(-1,0,1,0)+x_{4}(0,-1,0,1)$.
So $V$ is spanned by $(-1,0,1,0)$ and $(0,-1,0,1)$. It's eaisly
checked that these vectors are linearly indpendent. So $V$ is a basis
with $\dim V=2$.
\item $V=\mbox{span}\left\{ (3,-2,4),(-3,2,-4),(-6,4,-8)\right\} $. 
\[
\begin{bmatrix}\begin{array}{rrr}
3 & -3 & -6\\
-2 & 2 & 4\\
4 & -4 & -8
\end{array}\end{bmatrix}\rightsquigarrow\begin{bmatrix}\begin{array}{rrr}
1 & -1 & -2\\
0 & 0 & 0\\
0 & 0 & 0
\end{array}\end{bmatrix}
\]
So $(3,-2,4)$ (as is the other vectors) is a basis and$\dim V=1$.
\end{enumerate}

\end{document}
